% GENERATED FILE. Edit canonical lessons, then run npm run generate:books.
% canonical-source-hash: fnv1a64:654e04399273d4ec
% canonical-lessons: KA-C331-hindomme, KA-C331-andu, KA-C331-bahukala, KA-C331-pracina, KA-C331-yuga

\chapter{Long Ago}
\label{ch:ka-long-ago}
\hlchaptermodality{331}

\begin{chapteropening}
\textbf{By the end of this chapter:} \emph{I can say once before, on that day, a long time, ancient and an era.}

Complete the last lesson of chapter 331: I can say once before, on that day, a long time, ancient and an era.
\end{chapteropening}

\section[\texorpdfstring{\kn{ಹಿಂದೊಮ್ಮೆ} (hindomme)}{ಹಿಂದೊಮ್ಮೆ (hindomme)}]{\kn{ಹಿಂದೊಮ್ಮೆ} (hindomme) — once before}

\label{lesson:KA-C331-hindomme}



\begin{quote}
Before the new one: say the Kannada for the deceased, then the Kannada for a mishap.
\end{quote}

\subsection*{You'll want to know: \kn{ಹಿಂದೊಮ್ಮೆ}}
\textbf{\kn{ಹಿಂದೊಮ್ಮೆ}} — \emph{hindomme} — \textquotedblleft{}once before\textquotedblright{}.

\begin{quote}
\kn{ನಾನು} \kn{ಹಿಂದೊಮ್ಮೆ} \kn{ಈ} \kn{ಕೋಟೆಯನ್ನು} \kn{ನೋಡಿದ್ದೆ}.
\end{quote}

\begin{quote}
\emph{nānu hindomme ī kōṭeyannu nōḍidde.}
\end{quote}

\begin{quote}
I had seen this fort once before.
\end{quote}

\kn{ಹಿಂದೆ}, before, and \kn{ಒಮ್ಮೆ}, once.

\begin{cousinweb}
\textbf{\kn{ಹಿಂದೊಮ್ಮೆ}} (\emph{hindomme}) is a native Dravidian compound: \emph{hinde}, \textquotedblleft{}behind, in the past\textquotedblright{}, from \emph{pin} (DEDR 4205), and \kn{ಒಮ್ಮೆ}, \textquotedblleft{}once\textquotedblright{}.
\end{cousinweb}

\begin{grammarlens}[title={What you've built}]
One more word for telling what happened.
\end{grammarlens}

\subsection*{Your turn}
\begin{itemize}
\raggedright
  \item \emph{Say it:} \emph{hindomme}
  \item \emph{Say it:} \emph{hindomme}, once more
  \item \emph{Say it:} \emph{nānu hindomme ī kōṭeyannu nōḍidde.} — I had seen this fort once before.
  \item \emph{Recall:} say the Kannada for the deceased, then the Kannada for a mishap, then say \emph{hindomme} again
\end{itemize}

\begin{tcolorbox}[breakable,colback=teal!4,colframe=teal!35!black,title={Before you move on}]
What does \kn{ಹಿಂದೊಮ್ಮೆ} mean? (\textquotedblleft{}Once before\textquotedblright{}.) Say it once more.
\end{tcolorbox}
\section[\texorpdfstring{\kn{ಅಂದು} (andu)}{ಅಂದು (andu)}]{\kn{ಅಂದು} (andu) — on that day; then}

\label{lesson:KA-C331-andu}



\begin{quote}
Before the new one: say the Kannada for a mishap, then the Kannada for once before.
\end{quote}

\subsection*{You'll want to know: \kn{ಅಂದು}}
\textbf{\kn{ಅಂದು}} — \emph{andu} — \textquotedblleft{}on that day; then\textquotedblright{}.

\begin{quote}
\kn{ಅಂದು} \kn{ಊರಿನಲ್ಲಿ} \kn{ದೊಡ್ಡ} \kn{ಹಬ್ಬ} \kn{ಇತ್ತು}.
\end{quote}

\begin{quote}
\emph{andu ūrinalli doḍḍa habba ittu.}
\end{quote}

\begin{quote}
On that day there was a big festival in the village.
\end{quote}

\kn{ಇಂದು} is today, \kn{ಅಂದು} is that day. Stories often begin \kn{ಅಂದು} …

\begin{cousinweb}
\textbf{\kn{ಅಂದು}} (\emph{andu}) is native Dravidian: it is built on the pointing base \emph{a-}, \textquotedblleft{}that, yonder\textquotedblright{} (DEDR 1), as \kn{ಇಂದು}, \textquotedblleft{}today\textquotedblright{}, is built on \emph{i-}, \textquotedblleft{}this\textquotedblright{}.
\end{cousinweb}

\begin{grammarlens}[title={What you've built}]
One more word for telling what happened.
\end{grammarlens}

\subsection*{Your turn}
\begin{itemize}
\raggedright
  \item \emph{Say it:} \emph{andu}
  \item \emph{Say it:} \emph{andu}, once more
  \item \emph{Say it:} \emph{andu ūrinalli doḍḍa habba ittu.} — On that day there was a big festival in the village.
  \item \emph{Recall:} say the Kannada for a mishap, then the Kannada for once before, then say \emph{andu} again
\end{itemize}

\begin{tcolorbox}[breakable,colback=teal!4,colframe=teal!35!black,title={Before you move on}]
What does \kn{ಅಂದು} mean? (\textquotedblleft{}On that day; then\textquotedblright{}.) Say it once more.
\end{tcolorbox}
\section[\texorpdfstring{\kn{ಬಹುಕಾಲ} (bahukāla)}{ಬಹುಕಾಲ (bahukāla)}]{\kn{ಬಹುಕಾಲ} (bahukāla) — a long time}

\label{lesson:KA-C331-bahukala}



\begin{quote}
Before the new one: say the Kannada for once before, then the Kannada for on that day.
\end{quote}

\subsection*{You'll want to know: \kn{ಬಹುಕಾಲ}}
\textbf{\kn{ಬಹುಕಾಲ}} — \emph{bahukāla} — \textquotedblleft{}a long time\textquotedblright{}.

\begin{quote}
\kn{ಅಜ್ಜ} \kn{ಬಹುಕಾಲ} \kn{ಆ} \kn{ಹಳ್ಳಿಯಲ್ಲಿ} \kn{ಇದ್ದರು}.
\end{quote}

\begin{quote}
\emph{ajja bahukāla ā haḷḷiyalli iddaru.}
\end{quote}

\begin{quote}
Grandfather lived in that village for a long time.
\end{quote}

A formal phrase; in talk people say \kn{ತುಂಬಾ} \kn{ಸಮಯ}.

\begin{cousinweb}
\textbf{\kn{ಬಹುಕಾಲ}} (\emph{bahukāla}) is borrowed from Sanskrit, a compound of \emph{bahu}, \textquotedblleft{}much, many\textquotedblright{} (CDIAL 9187), and \emph{kāla}, \textquotedblleft{}time\textquotedblright{} (Monier-Williams).
\end{cousinweb}

\begin{grammarlens}[title={What you've built}]
One more word for telling what happened.
\end{grammarlens}

\subsection*{Your turn}
\begin{itemize}
\raggedright
  \item \emph{Say it:} \emph{bahukāla}
  \item \emph{Say it:} \emph{bahukāla}, once more
  \item \emph{Say it:} \emph{ajja bahukāla ā haḷḷiyalli iddaru.} — Grandfather lived in that village for a long time.
  \item \emph{Recall:} say the Kannada for once before, then the Kannada for on that day, then say \emph{bahukāla} again
\end{itemize}

\begin{tcolorbox}[breakable,colback=teal!4,colframe=teal!35!black,title={Before you move on}]
What does \kn{ಬಹುಕಾಲ} mean? (\textquotedblleft{}A long time\textquotedblright{}.) Say it once more.
\end{tcolorbox}
\section[\texorpdfstring{\kn{ಪ್ರಾಚೀನ} (prācīna)}{ಪ್ರಾಚೀನ (prācīna)}]{\kn{ಪ್ರಾಚೀನ} (prācīna) — ancient}

\label{lesson:KA-C331-pracina}



\begin{quote}
Before the new one: say the Kannada for on that day, then the Kannada for a long time.
\end{quote}

\subsection*{You'll want to know: \kn{ಪ್ರಾಚೀನ}}
\textbf{\kn{ಪ್ರಾಚೀನ}} — \emph{prācīna} — \textquotedblleft{}ancient\textquotedblright{}.

\begin{quote}
\kn{ಇದು} \kn{ಒಂದು} \kn{ಪ್ರಾಚೀನ} \kn{ದೇವಸ್ಥಾನ}.
\end{quote}

\begin{quote}
\emph{idu ondu prācīna dēvasthāna.}
\end{quote}

\begin{quote}
This is an ancient temple.
\end{quote}

It stands before a noun as it is: \kn{ಪ್ರಾಚೀನ} \kn{ಕಾಲ}, ancient times.

\begin{cousinweb}
\textbf{\kn{ಪ್ರಾಚೀನ}} (\emph{prācīna}) is borrowed from Sanskrit \emph{prācīna}, \textquotedblleft{}turned to the front, eastern\textquotedblright{} (CDIAL 8920), which also came to mean \textquotedblleft{}former, ancient\textquotedblright{}.
\end{cousinweb}

\begin{grammarlens}[title={What you've built}]
One more word for telling what happened.
\end{grammarlens}

\subsection*{Your turn}
\begin{itemize}
\raggedright
  \item \emph{Say it:} \emph{prācīna}
  \item \emph{Say it:} \emph{prācīna}, once more
  \item \emph{Say it:} \emph{idu ondu prācīna dēvasthāna.} — This is an ancient temple.
  \item \emph{Recall:} say the Kannada for on that day, then the Kannada for a long time, then say \emph{prācīna} again
\end{itemize}

\begin{tcolorbox}[breakable,colback=teal!4,colframe=teal!35!black,title={Before you move on}]
What does \kn{ಪ್ರಾಚೀನ} mean? (\textquotedblleft{}Ancient\textquotedblright{}.) Say it once more.
\end{tcolorbox}
\section[\texorpdfstring{\kn{ಯುಗ} (yuga)}{ಯುಗ (yuga)}]{\kn{ಯುಗ} (yuga) — an era; an age}

\label{lesson:KA-C331-yuga}



\begin{quote}
Before the new one: say the Kannada for a long time, then the Kannada for ancient.
\end{quote}

\subsection*{You'll want to know: \kn{ಯುಗ}}
\textbf{\kn{ಯುಗ}} — \emph{yuga} — \textquotedblleft{}an era; an age\textquotedblright{}.

\begin{quote}
\kn{ಇದು} \kn{ಕಂಪ್ಯೂಟರ್} \kn{ಯುಗ}.
\end{quote}

\begin{quote}
\emph{idu kampyūṭar yuga.}
\end{quote}

\begin{quote}
This is the computer age.
\end{quote}

Plural: \kn{ಯುಗಗಳು}, \emph{yugagaḷu}.

\kn{ಹೊಸ} \kn{ಯುಗ} is a new era. In talk \kn{ಯುಗ} can also mean a very long time.

\begin{cousinweb}
\textbf{\kn{ಯುಗ}} (\emph{yuga}) is borrowed from Sanskrit \emph{yuga}, \textquotedblleft{}a yoke, a pair\textquotedblright{} (CDIAL 10482), which also meant \textquotedblleft{}an age of the world\textquotedblright{} (Monier-Williams).
\end{cousinweb}

\begin{grammarlens}[title={What you've built}]
One more word for telling what happened.
\end{grammarlens}

\subsection*{Your turn}
\begin{itemize}
\raggedright
  \item \emph{Say it:} \emph{yuga}
  \item \emph{Say it:} \emph{yuga}, once more
  \item \emph{Say it:} \emph{idu kampyūṭar yuga.} — This is the computer age.
  \item \emph{Recall:} say the Kannada for a long time, then the Kannada for ancient, then say \emph{yuga} again
\end{itemize}

\begin{tcolorbox}[breakable,colback=teal!4,colframe=teal!35!black,title={Before you move on}]
What does \kn{ಯುಗ} mean? (\textquotedblleft{}An era; an age\textquotedblright{}.) Say it once more.
\end{tcolorbox}
